%! TEX root = part2.tex
% the main TeX file which is intended to compile, :VimtexReload after adjustment
% :h vimtex-tex-root
\documentclass[../main.tex]{subfiles}
% \includeonlyframes{current}
% \setbeameroption{hide notes} % Only slides
% \setbeameroption{show only notes} % Only notes
% \setbeameroption{show notes on second screen=right} % Both

\begin{document}

\begin{frame}[t,mycolor=digiPH_leaf,mytitle=standard,dark]
	Pondasi adalah elemen struktur bawah \textit{(sub-structure)} yang berfungsi menyalurkan beban dari struktur atas ke tanah.
	Fungsi utama:

	\begin{itemize}
		\item Menyalurkan beban vertikal
		\item Menahan gaya lateral
		\item Menjaga kestabilan bangunan
		\item Mencegah penurunan berlebih
	\end{itemize}
\end{frame}

\begin{frame}[t,mycolor=digiPH_ocean,mytitle=center,dark]
	\frametitle{Klasifikasi Pondasi}

	Pondasi dibedakan menjadi:

	\begin{itemize}
		\item Pondasi Dangkal
		\item Pondasi Dalam
	\end{itemize}

	Pemilihan ditentukan oleh:
	\begin{itemize}
		\item Kondisi tanah
		\item Kedalaman tanah keras
		\item Besar beban bangunan
		\item Jenis struktur
	\end{itemize}


\end{frame}


\begin{frame}[t,mycolor=digiPH_leaf,mytitle=standard,dark]
	\frametitle{Pengertian Pondasi Dangkal}

	Digunakan apabila:

	\begin{itemize}
		\item Lapisan tanah keras dekat permukaan
		\item Beban bangunan relatif kecil–menengah
		\item Kedalaman ± 1–3 meter
	\end{itemize}

	Lebih ekonomis dan sederhana dalam pelaksanaan.

\end{frame}

\end{document}
